\chapter{Ensembles, Bagging and Random Forests}
\companion{StatQuest: Random Forests Part 1, Building, Using and Evaluating}{\ytid{J4Wdy0Wc_xQ}}

Random forests were asked about in InfoEdge interviews in detail (Glassdoor): intuition, bootstrap sampling, random
feature selection, how each tree is trained, why a forest reduces variance compared with one tree, RF vs gradient
boosting, and feature importance. The OA asked whether RF is bagging or boosting and the types of ensembles. The
2026--27 talk even jotted ``Random forest \dots\ probabilistic''. This chapter answers all of that from first principles.

\section{Wisdom of crowds}
Ask one person to guess the number of sweets in a jar: they may be far off. Average the guesses of 500 people and the
answer is usually remarkably close. Two conditions make this work: each guess must be \emph{somewhat} informed (better
than random), and the errors must be \emph{different} from person to person (not all fooled the same way). Ensembles
apply this to models.

\section{Three families of ensembles}
\begin{itemize}
\item \textbf{Bagging} (bootstrap aggregating): train many models independently, each on a random resample of the data,
  and average them (or vote). Reduces \emph{variance}. Example: random forest.
\item \textbf{Boosting}: train models one after another, each focusing on the mistakes of the ensemble so far, and add
  them up. Reduces \emph{bias} (and variance somewhat). Examples: AdaBoost, gradient boosting, XGBoost, LightGBM, CatBoost.
\item \textbf{Stacking}: train several different models, then train a \emph{meta-model} that learns how to combine their
  predictions (using out-of-fold predictions to avoid leakage).
\end{itemize}
Also: \textbf{voting} (hard: majority class; soft: average probabilities) and \textbf{blending} (stacking with a single holdout
set instead of out-of-fold predictions).

\section{Why averaging reduces variance: the formula}
Suppose $B$ models each have prediction variance $\sigma^2$ and every pair has correlation $\rho$. The variance of their average is
\[ \Var\Big(\frac1B\sum_b \hat f_b\Big) = \rho\sigma^2 + \frac{1-\rho}{B}\sigma^2. \]
\begin{itemize}
\item If the models were independent ($\rho = 0$): $\sigma^2/B$. Ten models cut variance tenfold.
\item If identical ($\rho = 1$): $\sigma^2$. No gain at all.
\item As $B\to\infty$, the variance falls to $\rho\sigma^2$: \emph{correlation} sets the floor.
\end{itemize}
So an ensemble wants many, individually strong, \emph{decorrelated} models. Averaging does not change the bias (the
average of models with the same bias has that bias).

\begin{example}[: how much does averaging help?]
$\sigma^2 = 4$, $B = 10$. Independent: $0.4$. With $\rho = 0.3$: $1.2 + 0.7\cdot4/10 = 1.2 + 0.28 = 1.48$. With 1{,}000 models
and $\rho = 0.3$: $1.2 + 0.0028 \approx 1.2$. Beyond a point, more trees barely help; lower correlation does.
\end{example}

\section{Majority voting: the maths of the crowd}
If $B$ independent classifiers each have accuracy $p > 0.5$, majority vote accuracy rises with $B$ (Condorcet's jury theorem).
\begin{example}[: three voters]
Each 80\% accurate, independent. The vote is right if at least 2 of 3 are right:
$P = 0.8^3 + 3\cdot0.8^2\cdot0.2 = 0.512 + 0.384 = 0.896$. With five voters at 0.6: $P = \sum_{k=3}^5\binom5k0.6^k0.4^{5-k} =
0.3456 + 0.2592 + 0.0778 = 0.683$.
If each were only 40\% accurate, three voters give $0.4^3 + 3\cdot0.4^2\cdot0.6 = 0.352$: voting makes bad models worse.
\end{example}
Real models are correlated, so gains are smaller, which is again why random forests work hard to decorrelate trees.

\section{Bootstrap sampling}
A \textbf{bootstrap sample} draws $n$ rows \emph{with replacement} from $n$ rows. Some rows appear several times, some not at all.
The probability that a particular row is never chosen is $(1 - 1/n)^n \to e^{-1} \approx 0.368$. So each bootstrap sample contains
about \textbf{63.2\%} of the distinct rows; the other 36.8\% are \textbf{out-of-bag (OOB)} for that model.
\begin{example}[: OOB probability]
$n = 1000$: $(0.999)^{1000} = 0.3677$. Expected unique rows in a bootstrap of 100: $100(1 - 0.99^{100}) = 63.4$.
\end{example}

\section{Random forest = bagged trees + random feature subsets}
\begin{enumerate}
\item For $b = 1,\dots,B$ (say 500):
  \begin{enumerate}
  \item Draw a bootstrap sample of the rows.
  \item Grow a deep tree on it (often unpruned), but \textbf{at each split consider only a random subset of $m$ features}
    (default $m = \sqrt d$ for classification, $d/3$ or all features for regression in various libraries).
  \end{enumerate}
\item Predict by majority vote (or averaging class probabilities) for classification, or averaging for regression.
\end{enumerate}

\textbf{Why the random feature subset?} Suppose one feature (say ``preferred city match'') is very strong. Plain bagging would put it
at the root of nearly every tree, so the trees look alike and are highly correlated, and the $\rho\sigma^2$ floor stays high. Forcing
each split to choose from a random subset sometimes hides the dominant feature, so other features get to shape the tree. The trees
become less correlated; their average has lower variance. Each individual tree gets slightly worse (a bit more bias), but the
forest gets better.

\begin{example}[: the subset size]
64 features, classification: $m = \sqrt{64} = 8$ features per split. If one feature dominates, the probability it is
\emph{not} available at a given split is $1 - 8/64 = 87.5\%$.
\end{example}

\section{Out-of-bag error: free validation}
Each row is OOB for about a third of the trees. Predict each row using only the trees that did not see it, and compare with its
label. That OOB error is an almost unbiased estimate of test error, without a separate validation set. It is also how RF
permutation importance was originally computed.

\section{Feature importance in a forest}
\begin{itemize}
\item \textbf{Mean decrease in impurity (MDI)}: average over trees of the total impurity decrease from splits on each feature.
  Fast; biased toward high-cardinality/continuous features.
\item \textbf{Permutation importance} (mean decrease in accuracy): shuffle a feature's values in OOB or validation data and see how
  much accuracy drops. Less biased; slower; correlated features share credit.
\item \textbf{SHAP}: per-prediction attributions; TreeSHAP is exact for forests.
\end{itemize}

\section{Variants and practical notes}
\begin{itemize}
\item \textbf{Extra Trees} (extremely randomised trees): also pick thresholds at random (one per candidate feature) instead of
  the best one; usually no bootstrap. Even less correlation, faster training, slightly more bias.
\item \textbf{Isolation Forest}: random trees with random splits used for anomaly detection: anomalies are isolated in fewer splits
  (shorter average path length).
\item \textbf{Hyperparameters that matter}: \code{n\_estimators} (more is never worse for accuracy, only slower; stops improving
  around 200--1000), \code{max\_features} (the decorrelation knob), \code{min\_samples\_leaf}/\code{max\_depth} (tree complexity),
  \code{class\_weight} or balanced bootstrap for imbalance.
\item \textbf{Parallel}: trees are independent, so training is embarrassingly parallel (\code{n\_jobs=-1}).
\item \textbf{Weaknesses}: large memory and slower prediction than one tree; no extrapolation; probabilities can be poorly
  calibrated (votes cluster away from 0 and 1); less accurate than well-tuned boosting on many tabular problems.
\end{itemize}

\section{Random forest vs gradient boosting}
\begin{center}\small
\begin{tabularx}{\linewidth}{l X X}
\toprule
& Random forest & Gradient boosting\\
\midrule
Ensemble type & bagging (parallel, independent trees) & boosting (sequential, each tree fixes the previous)\\
Trees & deep, low bias, high variance & shallow (depth 3--8), high bias, low variance\\
Mainly reduces & variance & bias (and variance via shrinkage)\\
More trees & never overfits by itself (error plateaus) & can overfit; use early stopping\\
Tuning & forgiving, good defaults & sensitive (learning rate, depth, trees, regularisation)\\
Training & parallel across trees & sequential across trees (parallel within a tree)\\
Typical accuracy on tabular & very good & usually best\\
Noise / outliers in labels & robust & can chase noisy labels\\
\bottomrule
\end{tabularx}
\end{center}

\section{Stacking done right}
Level 0: several diverse models (logistic regression, RF, LightGBM, $k$-NN). Generate \emph{out-of-fold} predictions for each
training row (via $k$-fold CV) so every prediction is made by a model that did not see that row. Level 1: a simple model (often
logistic regression) learns the best combination from these predictions. Without out-of-fold predictions, the meta-model learns
to trust whichever base model overfit the most: a leakage bug.

\begin{practical}[: ensembles at InfoEdge]
\begin{itemize}
\item \textbf{Lead scoring / telecalling} (99acres): random forests are a robust first model on mixed user-behaviour features;
  their OOB score gives quick validation; permutation importance explains drivers to the sales team.
\item \textbf{Fraud / fake-recruiter detection}: Isolation Forest flags unusual posting behaviour (hundreds of jobs in minutes,
  odd salary ranges) without labels.
\item \textbf{Resdex ranking}: the ``M10'' ranking blocks on the Resdex slide are typically gradient-boosted learning-to-rank
  models; RF is the baseline they must beat.
\end{itemize}
\end{practical}

\stuck{StatQuest: Bootstrapping, Main Ideas}{\ytid{Xz0x-8-cgaQ}}
\section{Interview questions}

\begin{qa}{\pyq{OA: is Random Forest bagging or boosting?} Is a random forest bagging or boosting? Justify.}
\oneline{Bagging: its trees are trained independently and in parallel on bootstrap samples and averaged; it adds random feature
selection at each split to decorrelate the trees.}
\why Boosting would train trees sequentially, each one on reweighted data or residuals from the ensemble so far. RF trees never see
each other's errors. The ``random'' in random forest refers to two sources of randomness: bootstrap rows and random feature subsets
per split.
\pitfall Calling RF ``boosting because it combines weak learners''. RF's trees are deep (strong but high-variance) learners.
\end{qa}

\begin{qa}{\pyq{OA: types of ensembles} Name the types of ensemble methods, with an example and the error component each mainly
attacks.}
\oneline{Bagging (random forest; reduces variance), boosting (AdaBoost, GBM, XGBoost, LightGBM; reduces bias), and stacking
(meta-learner over diverse models; combines strengths), plus simple voting and blending.}
\why Bagging averages high-variance models; boosting builds a strong model from weak, high-bias ones by sequential correction;
stacking learns weights or interactions between models' predictions. Mixture of experts is a related idea where a gating
model chooses which expert to trust per input.
\end{qa}

\begin{qa}{\pyq{Glassdoor InfoEdge: Random Forest intuition} Explain random forests end to end: bootstrap sampling, random feature
selection, how each tree is trained and how predictions are made.}
\oneline{Build hundreds of deep trees, each on a bootstrap sample of the rows and each choosing every split from a random subset
of features, then average their predictions (or vote); averaging many decorrelated, low-bias trees gives a low-variance model.}
\why
\textbf{Step 1, rows.} For each tree, draw $n$ rows with replacement: about 63\% unique rows; the rest are out-of-bag for that tree.
\textbf{Step 2, features.} At every node, sample $m$ of the $d$ features (e.g. $\sqrt d$) and find the best split only among them.
\textbf{Step 3, grow.} Usually deep (low bias); controlled by \code{min\_samples\_leaf}. No pruning by default.
\textbf{Step 4, predict.} Classification: majority vote or mean of leaf class proportions (scikit-learn averages probabilities).
Regression: mean of leaf means.
\textbf{Why it works.} Deep trees have low bias but high variance; averaging cuts variance to about $\rho\sigma^2 + (1-\rho)\sigma^2/B$;
the two randomisations lower $\rho$.
\worked $d = 64$, 500 trees, $m = 8$: each split sees 8 random features; each tree sees about 632 of 1{,}000 rows.
\pushes \emph{``Is RF probabilistic?''} Its \code{predict\_proba} is the average of leaf class proportions across trees: a
frequency, not a calibrated probability. It tends to be under-confident near 0 and 1 (averaging pulls toward the middle); calibrate
with isotonic regression if probabilities are used for decisions.
\end{qa}

\begin{qa}{\pyq{Glassdoor InfoEdge} Why does a random forest reduce variance and overfitting compared with a single tree?}
\oneline{Because averaging $B$ models with variance $\sigma^2$ and pairwise correlation $\rho$ gives variance $\rho\sigma^2 +
(1-\rho)\sigma^2/B$, and the forest's bootstrap and feature subsampling keep $\rho$ low, while averaging leaves bias unchanged.}
\why A single deep tree fits the noise of its particular sample; another sample gives a very different tree. Errors that come from
sample-specific noise point in different directions for different trees and cancel in the average. Errors that come from a shared
systematic blind spot (bias) do not cancel. Deep trees have little bias, so the forest's remaining error is small.
\worked Single tree: 100\% train, 78\% test. 500-tree RF: 100\% train, 87\% test. Note RF training accuracy can still be 100\%:
it does not ``fit less'', it generalises better.
\pushes \emph{``Can a random forest overfit?''} Adding trees cannot increase overfitting (the average converges), but the
individual trees can overfit; with very noisy labels, limit \code{min\_samples\_leaf}. Also it can overfit through leakage, like any
model.
\end{qa}

\begin{qa}{\pyq{Glassdoor InfoEdge} Random forest vs gradient boosting: differences, and when would you choose each?}
\oneline{RF averages deep, independent trees to cut variance and is robust with little tuning; gradient boosting adds shallow trees
sequentially, each fitting the current errors, to cut bias, and usually wins on accuracy when tuned, at the cost of sensitivity to
hyperparameters and noise.}
\why See the comparison table. Choose RF for a fast, strong baseline, small or noisy datasets, or when you need OOB validation and
easy parallelism. Choose GBM (LightGBM/XGBoost) for maximum accuracy on large tabular data, custom losses (ranking, Poisson,
quantile), and when you can afford tuning and early stopping.
\pushes \emph{``Why does GBM overfit with more trees and RF not?''} In RF each new tree is an independent sample; the average
converges to a limit. In GBM each new tree fits the remaining residuals, which eventually are noise: it keeps reducing training loss
by fitting that noise.
\end{qa}

\begin{qa}{\pyq{Glassdoor InfoEdge: feature importance} How does a random forest compute feature importance, and what are its
pitfalls?}
\oneline{By default (MDI) it sums the impurity decreases from all splits on each feature, averaged over trees; permutation
importance (drop in OOB/validation score when a feature is shuffled) is the more reliable alternative.}
\why MDI is computed on training data and favours features with many split points (continuous, high-cardinality), even pure noise.
Permutation importance is computed on held-out data and reflects real predictive value, but correlated features hide each other.
Use both; drop-column importance (retrain without the feature) is the gold standard but costly.
\end{qa}

\begin{qa}{What is the probability that a given row is out-of-bag for one tree when $n = 1000$? Why does it approach 0.368?}
\oneline{$(1 - 1/1000)^{1000} \approx 0.368$, because $(1 - 1/n)^n \to e^{-1}$.}
\why Each of the $n$ draws misses the row with probability $1 - 1/n$; draws are independent. The limit is the definition of $e^{-1}$.
Hence a bootstrap contains about 63.2\% unique rows.
\end{qa}

\begin{qa}{Compute the variance of the average of 10 models with variance 4 when (a) independent and (b) pairwise correlation 0.5.}
\oneline{(a) 0.4; (b) $0.5\cdot4 + 0.5\cdot4/10 = 2.2$.}
\why Correlation dominates: even infinitely many models with $\rho = 0.5$ only halve the variance.
\end{qa}

\begin{qa}{Three independent classifiers with 80\% accuracy vote. What is the ensemble accuracy? What if each is 40\% accurate?}
\oneline{0.896 and 0.352: voting helps better-than-chance models and hurts worse-than-chance ones.}
\end{qa}

\begin{qa}{What is OOB error and how does it compare with cross-validation?}
\oneline{OOB error predicts each row using only the trees that did not train on it, giving a nearly free estimate of generalisation
error, comparable to cross-validation with about 1/3 holdout per tree.}
\why It is slightly pessimistic because each OOB prediction uses only about a third of the trees. It is valid only if rows are
independent: with grouped or time data, OOB leaks just like a random split.
\end{qa}

\begin{qa}{\deep Would bagging help a linear regression? A $k$-NN with $k = 1$? Why or why not?}
\oneline{Bagging helps unstable, high-variance learners (deep trees, 1-NN somewhat, neural nets) and barely helps stable ones like
linear regression, whose predictions hardly change across bootstrap samples.}
\why For OLS, the bagged average is very close to the full-data OLS fit: variance is already low and the models are highly
correlated ($\rho\approx1$). For 1-NN, bootstrapping changes which neighbour is nearest for some points, so it helps a bit, but 1-NN is
surprisingly stable to bootstrap (each point is present with probability 63\%). Subsampling features (random subspaces) helps $k$-NN
more.
\end{qa}

\begin{qa}{\deep Why does a random forest choose $\sqrt d$ features per split and not just 1, or all $d$?}
\oneline{It trades individual tree strength against correlation: $m = d$ gives strong but highly correlated trees (plain bagging), $m
= 1$ gives very decorrelated but weak, nearly random trees; $\sqrt d$ is an empirical sweet spot, and the right $m$ should be tuned.}
\why Breiman showed the forest's error bound depends on the ratio $\bar\rho/s^2$ (mean correlation over squared strength). Smaller $m$
lowers both $\bar\rho$ and strength $s$. With many noisy features and few informative ones, a small $m$ often misses the informative
features at a split, so larger $m$ (like $d/3$ or $0.5d$) is better.
\end{qa}

\begin{qa}{\deep A random forest outputs probability 0.9. Is it calibrated? How would you check and fix it?}
\oneline{Not necessarily; RF probabilities are vote fractions that tend to be pulled away from 0 and 1; check with a reliability
curve or Brier score on held-out data and fix with isotonic or Platt calibration.}
\why For an RF to output 1.0, every tree must agree; feature subsampling makes some trees always disagree, so extremes are rare
(under-confidence at the edges). Conversely, very deep trees with few rows per leaf give overconfident leaves. \code{CalibratedClassifierCV}
fits the calibration map with cross-validation.
\end{qa}

\begin{qa}{Explain stacking and how to avoid leakage in it.}
\oneline{Stacking trains a meta-model on the predictions of base models; to avoid leakage, the meta-model must be trained on
out-of-fold predictions, never on predictions the base models made for rows they trained on.}
\why If a base model overfits (100\% training accuracy), its in-sample predictions are perfect, and the meta-model learns to trust it
completely; on test data it fails. With $k$-fold, each row's base prediction comes from a model trained on the other folds. Test-time
base predictions come from models retrained on all data (or the average of fold models).
\end{qa}

\begin{qa}{Hard voting vs soft voting: compute both for five classifiers giving $P(\text{class }1) = 0.9, 0.6, 0.4, 0.45, 0.3$.}
\oneline{Hard voting: two vote 1, three vote 0, so class 0; soft voting: average 0.53, so class 1.}
\why Soft voting uses confidence: one very confident model can outweigh several lukewarm ones. It needs reasonably calibrated
probabilities.
\end{qa}

\begin{qa}{How does Isolation Forest detect anomalies?}
\oneline{It builds random trees with random features and random split points; anomalies, being few and different, get isolated in
very few splits, so a short average path length means a high anomaly score.}
\why A normal point in a dense region needs many random cuts to be separated from its neighbours. An isolated point (a recruiter
posting 400 jobs in an hour) is cut off quickly. Score $s = 2^{-\E[h(x)]/c(n)}$; close to 1 means anomaly. No labels needed,
linear time, works in high dimensions.
\end{qa}

\begin{qa}{\deep With 500 features of which only 2 are informative, how will a random forest behave and what would you do?}
\oneline{With $m = \sqrt{500}\approx 22$, most splits will not even see the two informative features, so many splits are on noise and
accuracy suffers; increase $m$, do feature selection first, or use boosting, which picks the best feature at each split.}
\why Probability a given informative feature is in a split's subset is $22/500 = 4.4\%$; both absent most of the time. Trees still
find them sometimes, but noise splits add variance. Remedies: raise \code{max\_features}, remove noise features (permutation or
Boruta), or use L1 logistic regression / GBM.
\end{qa}

\begin{qa}{Which hyperparameters of a random forest matter most and how would you tune them?}
\oneline{\code{max\_features} (decorrelation), \code{min\_samples\_leaf}/\code{max\_depth} (tree complexity), and enough
\code{n\_estimators}; tune the first two with CV or OOB, set trees high enough that OOB error has plateaued.}
\why \code{n\_estimators} is not a regularisation parameter: more trees only reduce noise in the average. \code{class\_weight=
"balanced\_subsample"} helps imbalance. \code{max\_samples} (bootstrap size) can speed training and add diversity.
\end{qa}

\begin{qa}{\deep Is a random forest a ``black box''? How would you explain one prediction to a recruiter?}
\oneline{It is hard to read globally but easy to explain locally: TreeSHAP gives each feature's contribution to a specific prediction,
and you can show the top contributing features in plain words.}
\why Example output for a Resdex candidate: base score 0.10; +0.25 because skills match 9 of 10 required; +0.12 because notice period
is 15 days; $-$0.05 because expected salary is above budget. This is the ``Explainability: why this candidate, in plain language''
idea from the PremiumX slide.
\end{qa}

\begin{qa}{What is the difference between bagging and pasting? Between random forest and Extra Trees?}
\oneline{Bagging samples rows with replacement, pasting without; Extra Trees additionally choose split thresholds at random and usually
use the whole dataset per tree, which further decorrelates trees and speeds training.}
\why Extra Trees trade a little bias for less variance and are often as accurate as RF and faster, especially with many features.
\end{qa}
